\section{Reproducibility Statement}\label{app:reproducibility}

All results in this paper are produced by deterministic scripts with fixed seeds.\footnote{Software attached as supplementary material in submission.} The 52-pair MLP lineage benchmark (Tables~\ref{tab:survival}, \ref{tab:baselines}) uses depth-16 MLPs (in\_dim=16, hidden=48) and is constructed from $2$ trained reference checkpoints, each paired with $5$ related-checkpoint kinds $\times 3$ derivative seeds $=30$ truly-related pairs (fine-tune, fine-tune to a new target, $\sigma_{\text{rel}} \in [0.01,0.15]$ Gaussian noise, $10$--$85\%$ magnitude pruning, and $16$--$256$-level uniform quantization), plus $16$ different-seed same-task pairs and $6$ distilled-student pairs as unrelated checkpoints. The expanded 159-pair set (Figure~\ref{fig:hero}c) uses $3$ depth-24 MLP reference models: each reference is paired with $25$ related checkpoints ($5$ transformation types $\times$ $5$ variations) and $28$ unrelated checkpoints ($15$ independent same-task $+$ $5$ independent different-task $+$ $5$ random-init $+$ $3$ distilled), yielding $3 \times 25 = 75$ related pairs and $3 \times 28 = 84$ unrelated pairs ($75 + 84 = 159$ total). Note that these 84 unrelated pairs share only 3 reference model roots; the effective sample size for null calibration is limited by the number of independent roots, not the number of pairs.

The GPT-2-Small-Lite benchmark (Table~\ref{tab:gpt2-survival}) uses 8 transformer roots with GPT-2 architecture ($L{=}6$ layers, $d{=}384$, $d_{\mathrm{ff}}{=}1536$, 6 heads, ${\sim}$35M parameters), trained on TinyStories~\citep{eldan2023tinystories} (50K samples, 3 epochs, batch size 8, lr$=$3e-4, cosine schedule, fp16, gradient checkpointing). Roots are allocated to calibration ($n{=}2$), development ($n{=}3$), and test ($n{=}3$) splits; Table~\ref{tab:gpt2-survival} reports results only on the 3 held-out test roots (21 descendant pairs, 3 distilled pairs, 21 cross-root pairs $=$ 45 total). Each root generates 7 descendants: 1 fine-tuned (1 epoch continued pretraining, lr$=$1e-4), 1 LoRA-merged (rank-8, $\alpha{=}16$, 1 epoch, lr$=$1e-4), 3 pruned (magnitude pruning at 30/50/70\% sparsity), and 2 quantized (uniform quantization at 256/64 levels, i.e., INT8/INT6). Additionally, 8 distilled students (one per root) are trained from scratch using knowledge distillation ($T{=}2.0$, $\alpha_{\mathrm{CE}}{=}0.5$, $\alpha_{\mathrm{KL}}{=}0.5$, 2 epochs). Quality metrics for distilled students include top-1 agreement, KL divergence, and perplexity ratio, computed on a held-out validation set.

Real-architecture pairing (Table~\ref{tab:transformer-family}) loads public HuggingFace checkpoints. Code, configuration files, baseline implementations (aligned Frobenius, singular-value distance, weight cosine, CKA, SVCCA, IPGuard regression analogue), the per-pair score JSONs, and the plotting scripts will be released upon acceptance.

\paragraph{Hardware and software.} All experiments were run on a single NVIDIA L4 GPU (24GB VRAM) with an AMD EPYC 7R13 CPU. Software: PyTorch, NumPy, SciPy, and Hugging Face Transformers. Global seeds (\texttt{torch.manual\_seed(0)}, \texttt{np.random.seed(0)}) are set at the start of each benchmark script; per-model seeds are derived deterministically from root and variant indices.

\paragraph{Data sources.} TinyStories~\citep{eldan2023tinystories} is loaded via Hugging Face Datasets (\texttt{roneneldan/TinyStories}). CIFAR-10~\citep{krizhevsky2009cifar} uses \texttt{torchvision.datasets}. Public checkpoints (GPT-2, LLaMA-2, Mistral, Qwen) are loaded from their official Hugging Face repositories.

\subsection{Latency Statistics}\label{app:latency}

Table~\ref{tab:gpt2-latency} reports per-method latency on the GPT-2 benchmark (20 pairs, GPU). Our method is 76$\times$ faster than Re-Basin+scale and 283$\times$ faster than SVD, while achieving identical AUROC under laundering.

\begin{table}[h]
\centering
\small
\begin{tabular}{@{}lrrrr@{}}
\toprule
Method & Mean & Std & Min & Max \\
\midrule
Centered Residual Sig.\ (ours) & \textbf{5.1} & 2.3 & 4.4 & 15.1 \\
Re-Basin+scale & 387.9 & 10.9 & 369.6 & 424.6 \\
Aligned Frobenius & 25.0 & 1.3 & 23.6 & 29.2 \\
Singular Value Distance & 1438.2 & 35.5 & 1417.7 & 1585.1 \\
Weight Cosine & 50.1 & 7.7 & 46.2 & 83.1 \\
\bottomrule
\end{tabular}
\caption{Per-method latency (ms) on GPT-2 benchmark. Our method achieves 76$\times$ speedup over Re-Basin+scale, the only other method maintaining AUROC$=$1.0 under laundering.}
\label{tab:gpt2-latency}
\end{table}

\subsection{Architecture-Aware Factorization}\label{app:factorization}

Table~\ref{tab:factorization} lists the architecture-aware residual branch products used throughout this work.

\begin{table*}[t]
\centering
\small
\begin{tabular}{@{}lll@{}}
\toprule
Architecture & Residual branch $F(x)$ & Product $M$ \\
\midrule
Transformer MLP / BasicBlock & $W_2 \sigma(W_1 x)$ & $W_2 W_1$ \\
Attention V/O path & $W_O \mathrm{Attn}(x) W_V x$ & $W_O W_V$ \\
Attention Q/K path & bilinear $x^\top W_Q W_K^\top x$ & $W_Q W_K^\top$ \\
Bottleneck ResNet & $W_3 \sigma(W_2 \sigma(W_1 x))$ & $W_3 W_2 W_1$ \\
SwiGLU MLP & $W_{\mathrm{down}} [\sigma(W_{\mathrm{gate}} x) \odot W_{\mathrm{up}} x]$ & $W_{\mathrm{down}} W_{\mathrm{up}}$ \\
\bottomrule
\end{tabular}
\caption{Architecture-aware residual branch products. Each $M$ composes the linear maps of one branch (dropping nonlinearities) so it maps the residual stream to itself. Factorization must match the architecture; incomplete products destroy the signal.}
\label{tab:factorization}
\end{table*}

On ImageNet-pretrained ResNet-50/101/152 (torchvision), we compare naive two-layer factorization ($W_3 W_1$, skipping the middle conv) against the correct three-layer product ($W_3 W_2 W_1$) on layer3 bottleneck blocks ($n{=}5/22/35$ blocks respectively). Naive $W_3 W_1$ achieves chance-level accuracy; the correct triple product recovers:
\begin{itemize}[nosep]
\item ResNet-50/layer3 ($n{=}5$): 100\%, AUC 1.000
\item ResNet-101/layer3 ($n{=}22$): 100\%, AUC 0.995
\item ResNet-152/layer3 ($n{=}35$): 91\%, AUC 0.964
\end{itemize}

\section{Verification Protocol Details}\label{app:verification}

\subsection{Gated Branch Score}

Branch similarity is reliable only when both branches exhibit strong trace concentration. We gate the cosine similarity:
\begin{equation}
G(\phi^A_\ell, \phi^B_m) = \langle \phi^A_\ell, \phi^B_m \rangle \cdot \min\!\left(\tfrac{s(M^A_\ell)}{\tau_s}, \tfrac{s(M^B_m)}{\tau_s}, 1\right)
\end{equation}
where $\tau_s$ is the minimum trace concentration across reference branches.

\subsection{Statistical Calibration}

We calibrate against a null distribution $\mathcal{N}_A = \{\mathcal{L}(A, U_j) : U_j \notin \mathcal{D}(A)\}$ of scores between reference $A$ and independently trained models. Given $n$ null scores, the empirical threshold $\tau = \max_j \mathcal{L}(A, U_j)$ yields a conformal p-value:
\begin{equation}
p(A,B) = \frac{|\{j : \mathcal{L}(A, U_j) \geq \mathcal{L}(A,B)\}| + 1}{n + 1}
\end{equation}
With $n$ null models, the minimum achievable p-value is $1/(n+1)$. We therefore report AUROC as the primary metric, which aggregates discrimination across all thresholds without requiring specific FPR claims.

\paragraph{Calibration power by benchmark.} The MLP benchmark has 84 unrelated pairs from 3 independent roots; the effective sample size for calibration is 3 roots, not 84 pairs, since descendants from the same root are not independent samples. The GPT-2-Small-Lite benchmark uses 8 roots total, with 3 held out for testing; each test root is compared against 7 independent roots, but the effective sample size for formal hypothesis testing remains limited. We report AUROC as a discrimination summary rather than threshold-based verification claims with specific false-positive rate guarantees. The LLaMA-2 case study includes 7 independently trained architectural clones (OpenLLaMA v1/v2, Amber, Baichuan v1/v2, InternLM, Yi), all with identical architecture (32 layers, $d{=}4096$, intermediate$=$11008). All 7 score $|\mathcal{L}| < 5 \times 10^{-5}$ against the LLaMA-2-7B reference, while descendants score $\mathcal{L} \geq 0.3$---a ${\sim}10{,}000\times$ gap that provides meaningful null calibration for this architecture family.

\begin{table}[h]
\centering
\small
\begin{tabular}{@{}lccc@{}}
\toprule
Transformation & Mean $\mathcal{L}$ & Min $\mathcal{L}$ & $> \max(\text{null})$ \\
\midrule
Fine-tune / Quant / Noise & 0.99 & 0.97 & 45/45 \\
Fine-tune (diff.\ target) & 0.94 & 0.84 & 15/15 \\
Pruning (10--85\%) & 0.81 & 0.58 & 15/15 \\
\midrule
Independent / Distilled & 0.08 & 0.01 & - \\
\bottomrule
\end{tabular}
\caption{Lineage score survival under post-training modification ($n{=}75$ related, $n{=}84$ unrelated). All related checkpoints exceed the maximum null score ($\max \mathcal{L}_{\text{null}} = 0.20$); AUROC$=$1.000.}
\label{tab:survival}
\end{table}

\section{Mechanism Validation}\label{app:mechanism}

\subsection{Initialization Ablation}\label{app:init-ablation}

Table~\ref{tab:init-ablation} shows pair accuracy before and after training across seven initialization schemes on depth-24 residual MLPs (3 seeds, 200 epochs). All schemes show chance-level accuracy (${\leq}2.1\%$) at initialization, rising to $93$--$100\%$ after training. Orthogonal initialization provides dynamical isometry at init yet shows zero correct pairs before training.

\begin{table}[h]
\centering
\small
\begin{tabular}{@{}lccc@{}}
\toprule
Init scheme & Untrained & Trained & AUROC \\
\midrule
Orthogonal       & $0.0\%$  & $100\%$ & $0.947$ \\
Kaiming-normal   & $2.1\%$  & $97\%$  & $0.981$ \\
Kaiming-uniform  & $2.1\%$  & $98.6\%$ & $0.98$ \\
Xavier-normal    & $2.1\%$  & $100\%$ & $0.990$ \\
Xavier-uniform   & $2.1\%$  & $100\%$ & $0.99$ \\
Uniform          & $2.1\%$  & $93.8\%$ & --- \\
Gaussian $\sigma{=}0.02$ & $2.1\%$ & $13\%$ & $0.671$ \\
\bottomrule
\end{tabular}
\caption{Initialization ablation on depth-24 residual MLPs. All schemes show chance-level accuracy before training. The Gaussian $\sigma{=}0.02$ case fails because blocks collapse to near-zero contribution.}
\label{tab:init-ablation}
\end{table}

\subsection{GPT-2 Scaling, Trace, and Jacobian Analysis}\label{app:gpt2-details}

Table~\ref{tab:jacobian-orthogonality} reports the normalized Jacobian deviation $\delta_J^{\mathrm{norm}} = \|J^\top J/\|J\|_F^2 - I/d\|_F$ across GPT-2 scales. This metric is scale-invariant: any scaled orthogonal matrix scores 0, and larger values indicate non-uniform singular values.

\begin{table}[h]
\centering
\small
\setlength{\tabcolsep}{3pt}
\begin{tabular}{@{}lcccr@{}}
\toprule
Model & $d$ & Pretrained & Rand-init & Ratio \\
\midrule
GPT-2-small  & 768  & 0.297 & 0.025 & $11.7\times$ \\
GPT-2-medium & 1024 & 0.242 & 0.026 & $9.4\times$ \\
GPT-2-large  & 1280 & 0.155 & 0.025 & $6.2\times$ \\
GPT-2-XL     & 1600 & 0.122 & 0.024 & $5.1\times$ \\
\bottomrule
\end{tabular}
\caption{Jacobian orthogonality ($\delta_J^{\mathrm{norm}}$) across GPT-2 scales. Pretrained models are $5$--$12\times$ less orthogonal than at random init, inconsistent with the hypothesis that trace concentration results from blocks approaching isometry.}
\label{tab:jacobian-orthogonality}
\end{table}

Figure~\ref{fig:gpt2-pairing} shows block pairing score matrices $s(i,j)$ for GPT-2 models from 124M to 1.5B parameters.

\begin{figure*}[t]
\centering
\includegraphics[width=\textwidth]{fig_gpt2_mlp_pairing.pdf}
\caption{Block pairing score matrices $s(i,j)$ for GPT-2 models from 124M to 1.5B parameters. Diagonal entries dominate, yielding 100\% pairing accuracy across all scales.}
\label{fig:gpt2-pairing}
\end{figure*}

Figure~\ref{fig:gpt2-trace} shows the trace values $\mathrm{tr}(W_{\mathrm{out}} W_{\mathrm{in}})$ per block across all four GPT-2 variants.

\begin{figure*}[t]
\centering
\includegraphics[width=0.95\textwidth]{fig_gpt2_trace_lollipop.pdf}
\caption{Distribution of trace values $\mathrm{tr}(W_{\mathrm{out}} W_{\mathrm{in}})$ across GPT-2 scales. Blue: negative trace (86\% average); orange: positive trace (14\% average). The strong skew toward negative values is consistent with the identity-alignment mechanistic account.}
\label{fig:gpt2-trace}
\end{figure*}

Figure~\ref{fig:gpt2-methods} compares block pairing accuracy across GPT-2 scales.

\begin{figure}[t]
\centering
\includegraphics[width=0.75\columnwidth]{fig_gpt2_methods_slope.pdf}
\caption{Block pairing accuracy as model size increases. Normalized trace pairing (blue) maintains 100\% projection-pair accuracy; Frobenius matching (orange) degrades from 67\% to 52\%; random baseline (gray) drops from 8\% to 2\% as the number of blocks increases.}
\label{fig:gpt2-methods}
\end{figure}

\subsection{Centering Ablation}\label{app:centering-ablation}

The centering step removes the identity-aligned component shared by all trained models. We verify this materially improves discrimination by directly comparing cosine similarity on uncentered branch products $\cos(\mathrm{vec}(M_\ell^A), \mathrm{vec}(M_\ell^B))$ versus centered residuals $\cos(\mathrm{vec}(R_\ell^A), \mathrm{vec}(R_\ell^B))$ where $R_\ell = M_\ell - \tfrac{\mathrm{tr}(M_\ell)}{d} I$.

Table~\ref{tab:centering-ablation} shows results on the GPT-2 benchmark (27 pairs from 3 test roots). Both methods achieve AUROC$=$1.0, but centering reduces unrelated-pair similarity by ${\sim}50\times$: from 0.019 (uncentered) to 0.0004 (centered). For independent roots specifically, uncentered similarity averages 0.027 while centered similarity averages 0.0015---an $18\times$ reduction.

\begin{table}[h]
\centering
\small
\setlength{\tabcolsep}{3pt}
\begin{tabular}{@{}lcccc@{}}
\toprule
Method & AUC & Gap-$Z$ & Unrel.\ $\mu$ & Margin \\
\midrule
Uncentered $\mathrm{vec}(M_\ell)$ & 1.0 & $+28.7$ & 0.019 & 0.84 \\
Centered $\mathrm{vec}(R_\ell)$ & 1.0 & $+29.2$ & \textbf{0.0004} & \textbf{0.86} \\
\bottomrule
\end{tabular}
\caption{Centering ablation on the GPT-2 benchmark. Both methods perfectly separate related from unrelated pairs, but centering reduces spurious similarity by ${\sim}50\times$.}
\label{tab:centering-ablation}
\end{table}

The identity energy fraction $|\mathrm{tr}(M)|^2 / \|M\|_F^2$ averages 9.4 across all layers, confirming that the identity component carries substantial matrix energy. Independent models share this generic identity-aligned structure from training dynamics; centering removes it, revealing the checkpoint-specific remainder that differs between unrelated models.

\section{Extended Benchmarks}\label{app:benchmarks}

\subsection{Laundering Experiment Gap-Z Scores}\label{app:laundering-gapz}

Table~\ref{tab:laundering-gapz} reports the same laundering experiment as Table~\ref{tab:laundering} but with Gap-$Z$ instead of AUROC. Gap-$Z$ measures the standardized separation between descendant and non-descendant score distributions; larger values indicate more robust margins.

\begin{table}[h]
\centering
\small
\setlength{\tabcolsep}{3pt}
\begin{tabular}{@{}llccccc@{}}
\toprule
& & \multicolumn{5}{c}{Gap-$Z$} \\
\cmidrule(lr){3-7}
Cond. & Bench & \textbf{Ours} & Re-Basin & Al.~Frob & SVD & W.~Cos \\
\midrule
none & MLP & $+53.0$ & $+48.2$ & $+72.0$ & $+7.9$ & $+76.3$ \\
P & MLP & $+53.0$ & $+48.2$ & $-0.1$ & $+7.9$ & $+2.1$ \\
P & GPT-2 & $+31.0$ & $+29.5$ & $+3.9$ & $+0.4$ & $+0.8$ \\
\midrule
D$_\mathrm{m}$ & MLP & $+53.0$ & $+48.2$ & $+72.0$ & $-5.2$ & $+76.3$ \\
D$_\mathrm{s}$ & MLP & $+53.0$ & $+48.2$ & $-4.8$ & $-5.2$ & $+76.3$ \\
PD & MLP & $+53.0$ & $+48.2$ & $-4.8$ & $-5.2$ & $+1.8$ \\
PDFT & MLP & $+49.1$ & $+45.0$ & $-4.8$ & $-5.2$ & $+1.8$ \\
\bottomrule
\end{tabular}
\caption{Gap-$Z$ under function-preserving laundering (same experiment as Table~\ref{tab:laundering}). Weight cosine collapses from $+76.3$ to $+2.1$ under permutation and $+1.8$ under PD---the margin is razor-thin despite AUROC$=$0.80--0.86. Our signature maintains Gap-$Z \approx +53$ across all conditions.}
\label{tab:laundering-gapz}
\end{table}

The key finding: weight cosine's Gap-$Z$ collapses from $+76.3$ (none) to $+2.1$ (P) and $+1.8$ (PD). Despite preserving AUROC$=$0.80--0.86, the descendant and non-descendant distributions nearly overlap. One additional perturbation would likely break the ranking entirely. Our signature maintains Gap-$Z \approx +53$ across all conditions, indicating stable separation regardless of laundering.

\subsection{Harder Regime: Layer Grafts and Linear Merges}\label{app:harder-bench}

We construct 40 harder pairs from depth-16 MLPs (same setup as Table~\ref{tab:baselines}): (i) \textbf{layer grafts} copy $K \in \{0,4,8,12,16\}$ of 16 blocks from reference into an independent donor (related if $K \ge 8$); (ii) \textbf{linear merges} compute $w \cdot \mathrm{ref} + (1-w) \cdot \mathrm{donor}$ with $w \in \{0, 0.25, 0.5, 0.75, 1.0\}$ (related if $w \ge 0.5$). We use 2 references and 2 donors per reference. AUROC measures binary related-vs-unrelated detection; Spearman correlation measures whether a method tracks partial overlap monotonically ($K/L$ or $w$) rather than just clearing a threshold. Table~\ref{tab:harder-bench} summarizes the 40-pair benchmark.

\begin{table}[h]
\centering\small
\setlength{\tabcolsep}{2pt}
\begin{tabular}{@{}lcccc@{}}
\toprule
& \multicolumn{2}{c}{AUROC} & \multicolumn{2}{c}{Spearman $\rho$} \\
\cmidrule(lr){2-3}\cmidrule(lr){4-5}
Method & Graft & Merge & Graft & Merge \\
\midrule
Ours  & 0.98 & 1.00 & +.96 & +.96 \\
Aligned Frob.\ & 1.00 & 1.00 & +.99 & +.99 \\
Weight cos.\ & 1.00 & 1.00 & +.98 & +.96 \\
SVD dist.\ & 1.00 & 0.45 & +.99 & +.23 \\
SVCCA & 1.00 & 1.00 & +.97 & +.99 \\
CKA & 0.81 & 0.78 & +.73 & +.75 \\
IPGuard & 0.05 & 0.45 & $-$.83 & +.22 \\
\bottomrule
\end{tabular}
\caption{Harder-regime benchmark on layer grafts and linear merges (40 pairs). Our method tracks partial overlap monotonically ($\rho \ge 0.96$). CKA degrades; IPGuard collapses on grafts; SVD fails merges.}
\label{tab:harder-bench}
\end{table}

\subsection{CIFAR-10 ResNet-18 Benchmark}\label{app:cifar-resnet}

We train two CIFAR-10 ResNet-18 references and three independents. Each reference yields 11 related checkpoints (noise $1{-}10\%$, pruning $20{-}80\%$, quantization $16{-}256$ levels). Results: AUROC$=$1.000; all related checkpoints exceed the maximum null score. Related checkpoint scores: $\mathcal{L} \in [0.695, 1.000]$; unrelated baseline: $\mathcal{L}_{\max} = 0.004$.

\subsection{GPT-2-Small-Lite Benchmark}\label{app:gpt2-benchmark}

Figure~\ref{fig:gpt2-benchmark} visualizes the GPT-2-Small-Lite benchmark results from Section~\ref{sec:exp-posttraining}. Panel (a) shows lineage score distributions by transformation type: descendants (quantized, LoRA, fine-tuned, pruned) score high while non-descendants (distilled students, independent roots) score near zero. The minimum descendant score (0.855) exceeds the maximum non-descendant score (0.004) by $200\times$. Panel (b) shows that distillation modestly improves teacher--student agreement (79\% top-1, +1.4 points over independent training) yet leaves the weight-lineage score near the independent-root range ($\mathcal{L} \approx 0.002$). This illustrates the conceptual distinction: behavioral similarity does not imply weight ancestry.

\begin{figure}[h]
\centering
\includegraphics[width=\columnwidth]{fig_gpt2_lineage_benchmark.pdf}
\caption{GPT-2-Small-Lite lineage benchmark (8 roots, 120 pairs across all splits; Table~\ref{tab:gpt2-survival} reports results on the 3 held-out test roots, 45 pairs). (a) Lineage scores by transformation type. Descendants (green) score high; non-descendants (red) score near zero. (b) Distillation modestly improves teacher--student agreement (79\% top-1, +1.4 points over independent training) while leaving the weight-lineage score near the independent-root range, confirming that behavioral similarity $\neq$ weight inheritance.}
\label{fig:gpt2-benchmark}
\end{figure}

\subsection{ROC for Lineage Verification}

Figure~\ref{fig:lineage-roc} shows that trace concentration alone cannot distinguish related from unrelated trained models.

\begin{figure}[h]
\centering
\includegraphics[width=0.85\columnwidth]{fig_lineage_three_panel.pdf}
\caption{ROC for lineage verification on depth-24 residual MLPs. The residual-signature score achieves AUROC$=$1.000; trace concentration alone fails (AUROC$=$0.417) because it cannot distinguish two trained residual models.}
\label{fig:lineage-roc}
\end{figure}

\subsection{Gradient-Based Suppression Attack}\label{app:attacks}

An adversary optimizes perturbation $\Delta$ to minimize lineage score subject to utility:
\begin{equation}
\min_{\Delta} \;\mathcal{L}_{\mathrm{cos}}(A, A + \Delta) \;+\; \lambda \cdot \|f_{A+\Delta}(X) - f_A(X)\|_2^2
\end{equation}

We run this attack on depth-24 MLPs ($d{=}64$). The null baseline for this configuration (score between the reference and an independent same-task model) is $\mathcal{L}_{\text{null}} = 0.084$.

\begin{table}[h]
\centering
\small
\begin{tabular}{@{}rccl@{}}
\toprule
$\lambda$ & Final $\mathcal{L}$ & Eval loss & Verdict \\
\midrule
$0$         & 0.053 & 39.5 & utility destroyed \\
$10^{-2}$   & 0.054 & 0.77 & +12\% loss, below null \\
$10^{-1}$   & 0.083 & 0.70 & +1.5\% loss, $\approx$ null \\
$\ge 1$     & 0.91+ & 0.69 & utility preserved \\
\bottomrule
\end{tabular}
\caption{Pareto frontier for suppression attack. The null baseline for this setup is $\mathcal{L}_{\text{null}} = 0.084$. Reaching the null threshold costs +1.5\% utility loss ($\lambda{=}10^{-1}$); driving the score reliably below null ($\lambda{=}10^{-2}$) costs +12\%.}
\label{tab:attack-pareto}
\end{table}

\section{Transfer Beyond Language Models}\label{app:beyond-lm}

The trace concentration phenomenon underlying our method is not specific to language models. Table~\ref{tab:vision-speech} shows that architecture-aware branch products recover projection pairings in vision and speech architectures as well.

\begin{table}[h]
\centering
\footnotesize
\setlength{\tabcolsep}{4pt}
\begin{tabular}{@{}llcc@{}}
\toprule
Model & $L$ & Pair Acc & AUC \\
\midrule
ViT-B/16 & 12 & \textbf{100\%} & 1.00 \\
Whisper (tiny/base/sm) & 4--12 & \textbf{100\%} & 0.85 \\
ResNet-50/101/152 & 5--35 & 91--100\% & 0.96 \\
\bottomrule
\end{tabular}
\caption{Trace concentration in vision and speech architectures. Pair accuracy measures correct within-block projection recovery via Hungarian matching on $s(i,j)$. Random-init baselines: $\le 9\%$.}
\label{tab:vision-speech}
\end{table}

For ViT~\citep{dosovitskiy2020vit}, we use the same transformer MLP factorization as for language models ($W_2 W_1$). For Whisper~\citep{radford2023whisper}, we extract branch products from both encoder and decoder MLP blocks. For ResNet~\citep{he2016resnet}, bottleneck blocks require a three-layer product ($W_3 W_2 W_1$); naive two-layer factorization ($W_3 W_1$, skipping the middle conv) yields chance-level accuracy. On ImageNet-pretrained ResNet-50/101/152 (torchvision), the correct triple product recovers 91--100\% of block pairings in layer3.

These results demonstrate that the underlying signal---trace concentration in residual branch products---is a general property of trained residual networks, not specific to language model architectures. This universality motivates building lineage signatures on the trace-concentrated structure. We focus the main evaluation on language models to match the ACL venue, but the method applies wherever residual blocks exist.
